% Compile with:  pdflatex model.tex   (twice, for the table of contents)
\documentclass[11pt,a4paper]{article}

\usepackage[T1]{fontenc}
\usepackage[utf8]{inputenc}
\usepackage[margin=25mm]{geometry}
\usepackage{amsmath}
\usepackage{amssymb}
\usepackage{booktabs}
\usepackage[hidelinks]{hyperref}

\newcommand{\vect}[1]{\mathbf{#1}}
\newcommand{\quat}[1]{\mathbf{#1}}
\newcommand{\dd}{\mathrm{d}}
\newcommand{\Rot}{\vect{C}}
\newcommand{\eh}{\hat{\vect{e}}_h}
\newcommand{\ec}{\hat{\vect{e}}_c}
\newcommand{\code}[1]{\texttt{#1}}

\title{Booster Landing Simulation:\\Mathematical Model as Implemented}
\author{Generated from the source of \code{booster-landing}}
\date{Revision of the working tree, \today}

\begin{document}
\maketitle

\begin{abstract}
This document states the equations actually implemented in the
\code{booster-landing} repository, for external review. It is written from the
source, not from intent: where a code comment disagrees with the code, the code
is described and the discrepancy is flagged. Sections~\ref{sec:vehicle}
to~\ref{sec:closedloop} cover the C\texttt{++} guidance and plant taken
unmodified from \code{external/pdg}; Section~\ref{sec:donor} covers the
reconstruction layer this project adds on top of the donor simulator, which is
where most of the modelling risk sits; Section~\ref{sec:projection} covers the
projection used for display. Section~\ref{sec:limits} lists the assumptions and
open questions that most deserve scrutiny.
\end{abstract}

\tableofcontents

\section{Scope and conventions}

\subsection{What is and is not modelled here}

Two independent simulators feed one telemetry format.

\begin{description}
\item[\code{external/pdg}] Convex-optimisation powered descent guidance and a
  6-DoF rigid-body plant, in C\texttt{++}. Thrust and gravity only: there is no
  atmosphere, no aerodynamics and no contact model. Used unmodified.
\item[\code{external/booster-lander-simulator}] A C11 core with a US76
  atmosphere, body and grid-fin aerodynamics and leg contact, which flies the
  full return from stage separation. Used unmodified, as an executable; its
  internal equations are documented in its own repository and are \emph{not}
  restated here. What is restated is every quantity this project reconstructs
  from its output, in Section~\ref{sec:donor}.
\end{description}

\begin{description}
\item[\code{bridge/rtls\_trajectory.py}] Written for this project, in Python: a
  planar round-Earth point-mass model with a US76 atmosphere, which flies the
  solo first stage from liftoff through engine cutoff, the boostback burn and
  the ballistic coast, down to the entry interface at $65$~km. Its equations are
  given in full in Section~\ref{sec:rtls}. Below the entry interface it hands
  the flight to the donor.
\end{description}

\code{external/} carries \textbf{one} edit, and it is the subject of
Section~\ref{sec:donorpatch}: a \code{custom} scenario in the donor's
\code{core/scenario.c}, plus the \code{vx0} field it needs, because no shipped
scenario could express a stage that is still flying horizontally when it arrives.
Every shipped scenario is bit-identical under it. Nothing else in
\code{external/} is touched.

The producers are never mixed into a single plant; each writes the same
telemetry file format independently, and a handover between two of them is a
state vector, not a splice of equations of motion.

\subsection{Frames and notation}

The inertial frame is right-handed with $\hat{\vect{z}}$ up and the landing
target at the origin. The body frame has $\hat{\vect{z}}_B$ along the vehicle's
long axis, pointing from the engine toward the nose.

Quaternions are scalar-first, $\quat{q}=(q_0,q_1,q_2,q_3)$, unit norm, and map
body to inertial: $\vect{x}_I=\Rot(\quat{q})\,\vect{x}_B$. The product
$\otimes$ is the Hamilton product. $\vect{g}=(0,0,-g_0)$ with
$g_0=9.80665\ \mathrm{m/s^2}$.

SI units throughout: metres, seconds, kilograms, newtons, radians (degrees only
where a parameter is exposed in degrees).

\section{Vehicle model}
\label{sec:vehicle}

\subsection{State and control}

The state has dimension 14 and the control dimension 3:
\begin{equation}
\vect{x}=\bigl(m,\ \vect{r}_I,\ \vect{v}_I,\ \quat{q}_{B\to I},\ \boldsymbol{\omega}_B\bigr),
\qquad
\vect{u}=\vect{T}_B .
\end{equation}
The control is the thrust vector expressed in the body frame; its direction
relative to $\hat{\vect{z}}_B$ is the gimbal deflection and its magnitude is the
throttle setting.

\subsection{Equations of motion}

\begin{align}
\dot m &= -\frac{\lVert\vect{u}\rVert}{I_{sp}\,g_0} \;=\; -\alpha\lVert\vect{u}\rVert,
  &\alpha &\equiv \frac{1}{I_{sp} g_0}, \label{eq:mdot}\\
\dot{\vect{r}} &= \vect{v}, \label{eq:rdot}\\
\dot{\vect{v}} &= \frac{\Rot(\quat{q})\,\vect{u}}{m} + \vect{g}, \label{eq:vdot}\\
\dot{\quat{q}} &= \tfrac{1}{2}\,\quat{q}\otimes(0,\boldsymbol{\omega}), \label{eq:qdot}\\
\vect{J}\dot{\boldsymbol{\omega}} &= \vect{r}_T\times\vect{u}
  - \boldsymbol{\omega}\times\vect{J}\boldsymbol{\omega}. \label{eq:wdot}
\end{align}

$\vect{J}=\mathrm{diag}(J_x,J_y,J_z)$ is constant --- the inertia is \emph{not}
updated as propellant is consumed --- and $\vect{r}_T$ is the thrust application
point in the body frame, nominally $(0,0,-\ell)$ with $\ell=12$~m, so
$\vect{r}$ is the centre of mass. Since $\vect{r}_T$ is on the roll axis, a
single gimballed engine produces no roll torque; roll is left uncontrolled
throughout.

The tilt from vertical follows from the quaternion alone,
\begin{equation}
\cos\theta_{\mathrm{tilt}} = 1-2\left(q_1^2+q_2^2\right).
\label{eq:tilt}
\end{equation}

\section{Guidance}
\label{sec:guidance}

\subsection{Successive convexification with free final time}

The 6-DoF problem is solved by SCvx following Szmuk and A\c{c}{\i}kme\c{s}e
(arXiv:1802.03827). Free final time is handled by time dilation: the trajectory
is posed on normalised time $\tau\in[0,1]$ with
\begin{equation}
\frac{\dd\vect{x}}{\dd\tau} = \sigma\, \vect{f}(\vect{x},\vect{u}),
\qquad \sigma \equiv t_f ,
\end{equation}
so the time of flight becomes a scalar decision variable.

Each iteration linearises about the current reference
$(\bar{\vect{x}},\bar{\vect{u}},\bar\sigma)$:
\begin{equation}
\frac{\dd\vect{x}}{\dd\tau} \approx
  \vect{A}(\tau)\vect{x} + \vect{B}(\tau)\vect{u} + \vect{S}(\tau)\sigma + \vect{z}(\tau),
\end{equation}
\begin{equation}
\vect{A}=\bar\sigma\frac{\partial \vect{f}}{\partial \vect{x}},\quad
\vect{B}=\bar\sigma\frac{\partial \vect{f}}{\partial \vect{u}},\quad
\vect{S}=\vect{f}(\bar{\vect{x}},\bar{\vect{u}}),\quad
\vect{z}=-\left(\vect{A}\bar{\vect{x}}+\vect{B}\bar{\vect{u}}\right).
\end{equation}
The Jacobians are analytic (\code{rocketJacobians}). Discretisation uses a
first-order hold on $\vect{u}$ over $K$ nodes with the state-transition matrices
integrated by RK4 (\code{rk4Substeps} sub-steps per interval), giving
\begin{equation}
\vect{x}_{k+1} = \bar{\vect{A}}_k\vect{x}_k
 + \bar{\vect{B}}^-_k\vect{u}_k + \bar{\vect{B}}^+_k\vect{u}_{k+1}
 + \bar{\vect{S}}_k\sigma + \bar{\vect{z}}_k + \boldsymbol{\nu}_k ,
\end{equation}
where $\boldsymbol{\nu}_k$ is virtual control, which keeps the linearised
subproblem feasible at every iteration.

\subsection{Subproblem objective}

With $\hat{\cdot}$ denoting scaled variables and $s_{x_i}$ the state scale
factors,
\begin{equation}
\min\;\; -\frac{m_N}{m_{\mathrm{wet}}}
 \;+\; w_\nu\left(\sum_{k}\sum_{i}\frac{\lvert\nu_{k,i}\rvert}{s_{x_i}}
 \;+\; \sum_{k}\sum_{s}\zeta_{k,s}\right),
\label{eq:objective}
\end{equation}
maximising terminal mass, i.e.\ minimising fuel, with an $\ell_1$ penalty on the
dynamic defect and on the state-triggered-constraint slacks $\zeta$.

\subsection{Path constraints}

All are imposed at every node and are second-order-cone representable.

\begin{align}
m_k &\ge m_{\mathrm{dry}} && \text{(propellant budget)} \\
\lVert\vect{u}_k\rVert &\le T_{\max} && \text{(upper thrust bound)} \\
\bar{\vect{n}}_k^{\mathsf{T}}\vect{u}_k &\ge T_{\min},
  \quad \bar{\vect{n}}_k=\frac{\bar{\vect{u}}_k}{\lVert\bar{\vect{u}}_k\rVert}
  && \text{(lower bound, linearised)} \\
\lVert\vect{u}_k\rVert &\le \frac{u_{k,z}}{\cos\delta_{\max}}
  && \text{(gimbal envelope)} \\
\bigl\lVert(q_1,q_2)_k\bigr\rVert &\le \sqrt{\tfrac{1}{2}\left(1-\cos\theta_{\max}\right)}
  && \text{(tilt, from \eqref{eq:tilt})} \\
\lVert\boldsymbol{\omega}_k\rVert &\le \omega_{\max}
  && \text{(body rate)} \\
\bigl\lVert \vect{r}_{xy,k}-\vect{r}_{f,xy}\bigr\rVert
  &\le \cot\gamma_{gs}\,\left(r_{z,k}-r_{f,z}\right)
  && \text{(glide slope, if } \gamma_{gs}>0)
\end{align}

The lower thrust bound is the one genuinely non-convex constraint; it is
linearised about the reference direction rather than losslessly relaxed as in
the 3-DoF case. This is a place where the reference-dependence of the feasible
set is worth a reviewer's attention.

\subsection{State-triggered constraints}

A trigger--constraint pair $\bigl(g,c\bigr)$ encodes
``$g(\vect{x},\vect{u})<0 \Rightarrow c(\vect{x},\vect{u})\le 0$'' through the
single continuous inequality
\begin{equation}
h \;=\; -\min\!\left(g,0\right)\cdot c \;\le\; 0,
\end{equation}
linearised in each subproblem, after Szmuk, Reynolds and A\c{c}{\i}kme\c{s}e
(2020). Compound triggers compose $g$ by $\max$ (AND) or $\min$ (OR).

The shipped instance limits angle of attack above a speed threshold. With
$\hat{\vect{z}}_B=\Rot(\quat{q})\vect{e}_3$,
\begin{align}
g &= v_{\mathrm{trig}} - \lVert\vect{v}\rVert, \\
c &= \vect{v}\cdot\hat{\vect{z}}_B + \cos\alpha_{\max}\lVert\vect{v}\rVert .
\end{align}
$c\le0$ is equivalent to $\angle(-\vect{v},\hat{\vect{z}}_B)\le\alpha_{\max}$,
which is the angle of attack for a body flying nose-first.

\subsection{Trust region and acceptance}

The trust region is \emph{hard}, imposed as a cone constraint on the scaled
deviation at every node,
\begin{equation}
\left\lVert \bigl(\hat{\vect{x}}_k-\hat{\bar{\vect{x}}}_k;\ \hat{\vect{u}}_k-\hat{\bar{\vect{u}}}_k\bigr)\right\rVert \le r,
\qquad
\bigl\lvert \hat\sigma-\hat{\bar\sigma}\bigr\rvert \le r .
\end{equation}

\emph{Discrepancy flagged for review:} the header comment in
\code{include/pdg/scvx.hpp} describes ``soft quadratic trust regions''. The
implementation in \code{src/scvx.cpp} uses the hard cone above. The document
follows the code.

Steps are accepted or rejected on the ratio of actual to predicted penalty
reduction,
\begin{equation}
\rho = \frac{J(\bar{\vect{x}}) - J_{\mathrm{nl}}(\vect{x}^\star)}
            {J(\bar{\vect{x}}) - J_{\mathrm{lin}}(\vect{x}^\star)},
\end{equation}
accepting when $\rho\ge\rho_0$, shrinking the radius by $c_-$ when
$\rho<\rho_1$, growing it by $c_+$ when $\rho\ge\rho_2$, all clipped to
$[r_{\min},r_{\max}]$ (Mao, Szmuk and A\c{c}{\i}kme\c{s}e). The nonlinear
penalty $J_{\mathrm{nl}}$ is evaluated by single-shooting the accepted control
through the true dynamics, so the reported defect is a genuine integration
residual.

Convergence is declared when the scaled state change, the virtual control and
the predicted reduction all fall below their tolerances.

\subsection{The 3-DoF case, for reference}

A separate solver (\code{pdg::lcvx}) treats the translation-only minimum-fuel
problem by lossless convexification after A\c{c}{\i}kme\c{s}e and Ploen (2007):
the slack $\varsigma$ replaces $\lVert\vect{u}\rVert$ with
$\lVert\vect{u}\rVert\le\varsigma$, $T_{\min}\le\varsigma\le T_{\max}$, and the
relaxation is tight at the optimum. Free final time is a golden-section search
over $t_f$. This path is not used by the pipeline documented here and is
mentioned only for completeness.

\section{Plant and closed loop}
\label{sec:closedloop}

\subsection{Integration and control sampling}

The plant integrates \eqref{eq:mdot}--\eqref{eq:wdot} with fixed-step RK4. The
controller is sampled once per step at the step start and held constant across
it, which is a zero-order hold at $1/\Delta t$~Hz. The flight log is written from
inside that sampling, so each logged row is a real controller tick rather than a
re-integration.

Integration stops at the first step with $r_z\le0$; touchdown quantities are
interpolated between the bracketing states. Note that $\vect{r}$ is the centre of
mass and the model carries no vehicle height, so ``touchdown'' is the centre of
mass reaching ground level.

\subsection{Plant dispersions}

The plant is perturbed while the guidance and the autopilot are not told:

\begin{description}
\item[Thrust scale] $s_T$ multiplies the commanded vector.
\item[Specific impulse] $I_{sp}\leftarrow s_{I}\,I_{sp}$ in the plant only.
\item[Engine misalignment] a fixed cant $a$ about an axis in the body $x$--$y$
  plane at azimuth $\psi$, applied by Rodrigues' formula with
  $\hat{\vect{k}}=(-\sin\psi,\cos\psi,0)$:
  \begin{equation}
  \vect{R} = \vect{I}\cos a + [\hat{\vect{k}}]_\times \sin a
           + \hat{\vect{k}}\hat{\vect{k}}^{\mathsf{T}}(1-\cos a).
  \end{equation}
\item[Actuator lag] a first-order lag with time constant $\tau_a$ toward the
  mid-step target, applied exactly over the step:
  \begin{equation}
  \vect{u} \leftarrow \vect{u} + \left(1-e^{-\Delta t/\tau_a}\right)
    \left(\vect{u}_{\mathrm{tgt}}-\vect{u}\right).
  \end{equation}
\end{description}

The dispersed command is finally clamped to $\lVert\vect{u}\rVert\le T_{\max}$.

\subsection{Tracking autopilot}

The guidance solution is computed once; a PD autopilot flies it. This split is
deliberate --- an SCvx solve takes of order half a second and cannot sit inside a
control loop.

Let $(\vect{x}_r,\vect{u}_{r,B})$ be the reference interpolated to the current
time. The translation loop forms a saturated acceleration correction,
\begin{equation}
\vect{a}_c = \mathrm{sat}_{a_{\max}}\!\left[
  k_{p}\left(\vect{r}_r-\vect{r}\right) + k_{d}\left(\vect{v}_r-\vect{v}\right)\right],
\end{equation}
and a desired inertial force
\begin{equation}
\vect{f}_{\mathrm{des}} = \Rot(\quat{q}_r)\vect{u}_{r,B} + m\,\vect{a}_c,
\qquad
\lVert\vect{f}_{\mathrm{des}}\rVert \in [T_{\min},T_{\max}].
\end{equation}
Past the end of the reference the feedforward becomes a hover term
$\vect{f}_r=-m\vect{g}$.

The desired attitude is the reference attitude composed with the small rotation
carrying $\hat{\vect{f}}_r$ onto $\hat{\vect{f}}_{\mathrm{des}}$, using the
small-angle quaternion built from the cross product:
\begin{equation}
\vect{a}_{\mathrm{rot}}=\hat{\vect{f}}_r\times\hat{\vect{f}}_{\mathrm{des}},
\qquad
\dd\quat{q}=\mathrm{normalise}\bigl(1,\tfrac{1}{2}\vect{a}_{\mathrm{rot}}\bigr),
\qquad
\quat{q}_{\mathrm{des}}=\dd\quat{q}\otimes\quat{q}_r .
\end{equation}
The body-frame attitude error is $\quat{q}_e=\quat{q}^{-1}\otimes\quat{q}_{\mathrm{des}}$,
reduced to the two controllable components with a sign convention that takes the
short way round,
\begin{equation}
e_x = 2\,\mathrm{sgn}(q_{e,0})\,q_{e,1},\qquad
e_y = 2\,\mathrm{sgn}(q_{e,0})\,q_{e,2}.
\end{equation}
The PD attitude loop produces angular accelerations
$a_x = k_{a}e_x + k_{d,a}(\omega_{r,x}-\omega_x)$ and likewise $a_y$, converted
to a gimbal increment through the moment arm $L=\lvert r_{T,z}\rvert$:
\begin{equation}
u_x = u_{r,x}\frac{\lVert\vect{f}_{\mathrm{des}}\rVert}{\lVert\vect{f}_r\rVert} - \frac{J_y a_y}{L},
\qquad
u_y = u_{r,y}\frac{\lVert\vect{f}_{\mathrm{des}}\rVert}{\lVert\vect{f}_r\rVert} + \frac{J_x a_x}{L}.
\end{equation}
The lateral component is limited to half the commanded magnitude, and
\begin{equation}
u_z=\sqrt{\max\!\left(\lVert\vect{f}_{\mathrm{des}}\rVert^2-u_x^2-u_y^2,\;
  \tfrac{1}{4}\lVert\vect{f}_{\mathrm{des}}\rVert^2\right)} .
\end{equation}

\emph{For review:} this lateral limit of $\tfrac{1}{2}\lVert\vect{f}\rVert$
corresponds to a $30^\circ$ gimbal, far outside the $\delta_{\max}=12^\circ$ the
guidance was solved against. The autopilot therefore may command deflections the
guidance never considered feasible. It is a plant-side saturation only, but it
means the closed-loop trajectory is not guaranteed to respect the guidance
gimbal envelope.

\subsection{Feasibility bound reported before solving}

The bridge prints an energy screen so an impossible case is recognisable before
the solver is blamed. With no atmosphere, the entire deceleration must come from
the engine:
\begin{equation}
\Delta v_{\mathrm{avail}} = I_{sp} g_0 \ln\frac{m_{\mathrm{wet}}}{m_{\mathrm{dry}}},
\qquad
v_{\mathrm{arr}} = \sqrt{\lVert\vect{v}_0\rVert^2 + 2 g_0 h_0},
\qquad
\frac{T}{W} = \frac{T_{\max}}{m_{\mathrm{wet}} g_0}.
\end{equation}
$v_{\mathrm{arr}}$ is the speed an unpowered vehicle would reach at the ground
and is a strict lower bound on the required $\Delta v$; gravity losses during the
burn add to it. The case is reported as unachievable when
$v_{\mathrm{arr}}>\Delta v_{\mathrm{avail}}$.

\section{Donor reconstruction layer}
\label{sec:donor}

This section is the part of the model written for this project, and it carries
the most modelling risk. The donor core emits a text trajectory log at one line
per $250$ integration steps, i.e.\ $2$~Hz, carrying
\begin{equation}
\left\{t,\ h,\ v_z,\ \text{throttle},\ \lvert\theta_{\mathrm{tilt}}\rvert,\
 \ell,\ v_{\mathrm{rad}},\ \bar q,\ \omega_\perp,\ w,\ m,\ \text{phase}\right\},
\end{equation}
where $h$ is the height of the vehicle base, $\ell=\sqrt{r_x^2+r_y^2}$ is the
horizontal \emph{distance} from the pad, $v_{\mathrm{rad}}$ its rate, $\bar q$
dynamic pressure, and $\omega_\perp$ the magnitude of the perpendicular body
rate. Three quantities the viewer needs are absent and must be reconstructed.

\subsection{Thrust from mass flow}

The logged throttle fraction does not say how many engines are lit --- three for
the entry burn, one to land --- so thrust is taken from the mass flow instead:
\begin{equation}
T = \max\!\left(-\frac{\dd m}{\dd t},\,0\right) g_0\, I_{sp}(h),
\qquad
I_{sp}(h) = I_{sp}^{\mathrm{vac}}
 + \left(I_{sp}^{\mathrm{sl}}-I_{sp}^{\mathrm{vac}}\right) e^{-h/H},
\end{equation}
with scale height $H=8500$~m, and $T$ forced to zero wherever the logged
throttle is zero. The exponential is a crude stand-in for the ambient pressure
ratio; it affects only the $282$ to $311$~s blend, a $10\%$ effect on magnitude.

\subsection{Signing the range}

$\ell$ is a distance and cannot be negative, so a stage passing over the pad and
out the other side appears to reflect off $\ell=0$. The sign is recovered by
detecting the cusp this leaves. Between consecutive samples a crossing is
declared when all three hold:
\begin{equation}
v_{\mathrm{rad},i-1}\,v_{\mathrm{rad},i} < 0,
\qquad
\min\bigl(\lvert v_{\mathrm{rad},i-1}\rvert,\lvert v_{\mathrm{rad},i}\rvert\bigr) > v_\epsilon,
\qquad
\min\bigl(\ell_{i-1},\ell_i\bigr) < \max_i\lvert v_{\mathrm{rad}}\rvert\,\Delta t ,
\end{equation}
with $v_\epsilon=0.5$~m/s, after which the sign is inverted for all subsequent
samples. A genuine radial turnaround away from the pad has $v_{\mathrm{rad}}$
decaying smoothly through zero and fails the last two tests.

The unfolding is applied \emph{before} resampling, because the signed coordinate
passes smoothly through zero and is therefore safe to interpolate linearly,
while the folded one has a corner there.

This assumes the horizontal motion stays in one vertical plane. The log carries
no azimuth, so the assumption cannot be checked from the data.

\subsection{Lean direction}

The log carries the tilt magnitude but never which side it leans. Thrust acts
along the body axis, so a stage accelerating toward $+x$ is leaning that way:
\begin{equation}
\mathrm{sign}(\theta) = \mathrm{sign}\!\left(\overline{\dot v_{\mathrm{rad}}}\right)
\quad\text{while } T>0.05\,T_{\max},
\end{equation}
with the acceleration smoothed by a 5-sample boxcar and the last decision held
while the engine is off. The commanded lean angle is then
\begin{equation}
\theta_{\mathrm{cmd}} = \lvert\theta_{\mathrm{tilt}}\rvert \cdot \mathrm{sign}(\theta).
\end{equation}

\subsection{Attitude dynamics of the reconstruction}

$\theta_{\mathrm{cmd}}$ is a step function wherever the inferred sign changes,
and a stage cannot step. Drawing it directly produced $347^\circ$/s slews,
roughly twelve times anything the vehicle achieves. The command is therefore
flown through the control authority the vehicle actually has.

\subsubsection{Pitch inertia}

Following the donor's own \code{mass\_props}: the dry stage as a thin-walled
shell and the propellant as a solid column, both carried to the combined centre
of mass by the parallel-axis theorem. With propellant mass
$m_p=\max(m-m_{\mathrm{dry}},0)$, column cross-section $A$, mixture density
$\rho$, column radius $r_c^2=A/\pi$:
\begin{gather}
h_c = \frac{m_p}{\rho A},\qquad
z_c = z_{\mathrm{base}} + \tfrac{1}{2}h_c,\qquad
z_{\mathrm{com}} = \frac{m_{\mathrm{dry}}z_{\mathrm{dry}} + m_p z_c}{m},\\[2pt]
I_{\mathrm{dry}} = \frac{m_{\mathrm{dry}}\left(6R^2+L^2\right)}{12},
\qquad
I_{c} = \frac{m_p\left(3r_c^2+h_c^2\right)}{12},\\[2pt]
I_{\mathrm{pitch}} = I_{\mathrm{dry}} + m_{\mathrm{dry}}\left(z_{\mathrm{dry}}-z_{\mathrm{com}}\right)^2
 + I_{c} + m_p\left(z_c-z_{\mathrm{com}}\right)^2 .
\end{gather}
The log does not split LOX from RP-1, so a single column at the mixture density
for the nominal $O/F$ is used. Computed range over the \code{entry} profile:
$I_{\mathrm{pitch}}=5.1$ to $6.1\times10^6\ \mathrm{kg\,m^2}$, with
$z_{\mathrm{com}}$ moving from $11.6$ down to $7.4$~m as the tanks empty.

\subsubsection{Control authority}

Three torque sources, each dominant in a different phase:
\begin{align}
M_{\mathrm{gimbal}} &= T \sin\delta_{\max}\; z_{\mathrm{com}}, \\
M_{\mathrm{fin}} &= n_f\, \bar q\, S_f\, C_{N\alpha}\, \delta_f\,
  \left(z_f - z_{\mathrm{com}}\right), \\
M_{\mathrm{rcs}} &= n_r\, F_r \left(z_r - z_{\mathrm{com}}\right),
\end{align}
and the peak angular acceleration available is
\begin{equation}
\alpha_{\max} = \frac{M_{\mathrm{gimbal}}+M_{\mathrm{fin}}+M_{\mathrm{rcs}}}{I_{\mathrm{pitch}}} .
\label{eq:alphamax}
\end{equation}
Over the \code{entry} profile the peaks are $25.4$, $99.2$ and
$0.32\ \mathrm{deg/s^2}$ respectively: the fins fly the unpowered aero descent,
the gimbal the burns, and the RCS is all that remains in near-vacuum.

A slew of $\theta$ flown bang-bang --- accelerate half way, decelerate the rest
--- then takes
\begin{equation}
t_{\mathrm{slew}} = 2\sqrt{\frac{\theta}{\alpha_{\max}}},
\qquad
\omega_{\mathrm{peak}} = \sqrt{\theta\,\alpha_{\max}} .
\label{eq:slewtime}
\end{equation}

\subsubsection{Slew limiter}

The command is tracked by a time-optimal double integrator. At each step the
target rate is the fastest that can still be brought to rest exactly on target,
\begin{equation}
\omega_{\mathrm{des}} = \mathrm{sgn}(e)\sqrt{2\alpha_{\max}\lvert e\rvert},
\qquad e = \theta_{\mathrm{cmd}} - \theta,
\end{equation}
clipped to the logged rate bound, with the rate change limited by the available
acceleration:
\begin{equation}
\omega \leftarrow \mathrm{clip}\Bigl(
  \omega + \mathrm{clip}\bigl(\omega_{\mathrm{des}}-\omega,\ \pm\alpha_{\max}\Delta t\bigr),
  \ \pm\omega_\perp\Bigr),
\qquad
\theta \leftarrow \theta + \omega\,\Delta t .
\end{equation}

The cap at the logged $\omega_\perp$ is not a modelling choice: the in-plane
component of an angular rate cannot exceed the magnitude of that rate, so it is
an exact bound, and it keeps the reconstruction honest wherever
\eqref{eq:alphamax} is generous.

\paragraph{Validation.} The reconstructed peak rate sits just below the rate the
donor itself recorded, across every scenario:

\begin{center}
\begin{tabular}{lrr}
\toprule
scenario & reconstructed & logged $\omega_\perp$ \\
\midrule
\code{entry}       & $20.2^\circ$/s & $28.6^\circ$/s \\
\code{terminal}    & $8.4^\circ$/s  & $9.7^\circ$/s  \\
\code{aero\_offset}& $11.0^\circ$/s & $12.6^\circ$/s \\
\code{chaos}       & $11.9^\circ$/s & $12.0^\circ$/s \\
\bottomrule
\end{tabular}
\end{center}

\subsection{Synthesised attitude}

The reconstruction is a pure pitch about the inertial $y$ axis by $\theta$,
which takes $\hat{\vect{z}}_B$ toward $+x$:
\begin{equation}
\quat{q} = \left(\cos\tfrac{\theta}{2},\,0,\,\sin\tfrac{\theta}{2},\,0\right),
\qquad
\boldsymbol{\omega} = \left(0,\ \dot\theta,\ 0\right).
\end{equation}
The gimbal is written as zero, $\vect{u}=(0,0,T)$, because the log omits it.
Consequently the gimbal trace is identically zero for donor runs and the drawn
thrust is always along the body axis.

\subsection{Resampling}

The $2$~Hz log is interpolated linearly onto a uniform $20$~Hz grid for playback
only; the phase label is held rather than blended. This adds no information and
both rates are recorded in the metadata. All reconstruction above is performed on
the interpolated grid except the range unfolding, which precedes it.

\section{RTLS trajectory model}
\label{sec:rtls}

This section covers \code{bridge/rtls\_trajectory.py}, which flies the solo
first-stage return-to-launch-site profile from liftoff to the entry interface at
$h_{\mathrm{EI}}=65$~km on the way down. Below that the donor core takes the
flight, so Section~\ref{sec:donor} applies from there. The two are never mixed:
the handover is a state vector and nothing else. Why $65$ and not the $70$~km the
specification names is a property of the donor's guidance rather than of the
physics, and is set out in Section~\ref{subsec:handover}.

\subsection{Why the split is at 70\,km}

The boostback burn and the coast are a vacuum problem. Against US76 at
$1579\ \mathrm{m/s}$ the dynamic pressure is $386$~Pa at $60$~km, $23$~Pa at
$80$~km and $0.70$~Pa at $100$~km. On an $80$~t stage with the engines-first
$C_D S$ of Section~\ref{subsec:cds} the last of these is a deceleration of
$8\times10^{-5}\ \mathrm{m/s^2}$, four parts in $10^{5}$ of the thrust
acceleration. Thrust-and-gravity dynamics are therefore valid across the
boostback and the coast, and the atmosphere only has to be right on the way back
down --- which is exactly the donor's competence.

\subsection{Frame and equations of motion}

Planar, Earth-centred and \emph{non-rotating}, with the pad at
$(0,R_\oplus)$, $R_\oplus=6371$~km. The state is
$\vect{y}=(X,Z,V_X,V_Z,m)$ and gravity is central rather than flat, with $\mu$
pinned so that surface gravity is exactly $g_0$:
\begin{equation}
\mu = g_0 R_\oplus^2,
\qquad
\vect{a}_g = -\frac{\mu}{\lVert\vect{r}\rVert^3}\,\vect{r} .
\end{equation}
Altitude, ground-track range and the local frame are
\begin{equation}
h = \lVert\vect{r}\rVert - R_\oplus,
\qquad
d = R_\oplus \arctan\!\frac{X}{Z},
\qquad
\hat{\vect{u}} = \frac{\vect{r}}{\lVert\vect{r}\rVert},
\qquad
\hat{\vect{e}} = \frac{(Z,-X)}{\lVert\vect{r}\rVert} ,
\end{equation}
so that $\hat{\vect{u}}$ is local up and $\hat{\vect{e}}$ points downrange.
Curvature matters here: over $70$~km of ground track the surface falls $384$~m,
which is negligible against $130$~km of altitude but not against a landing.

The equations of motion are
\begin{equation}
\dot{\vect{v}} = \vect{a}_g
  + \frac{n_e\,\tau\,T_e(p)}{m}\,\hat{\vect{u}}_T
  - \frac{q S C_D(M)}{m}\,\frac{\vect{v}}{\lVert\vect{v}\rVert},
\qquad
\dot m = -n_e\,\tau\,\dot m_e ,
\end{equation}
with $q=\tfrac12\rho\lVert\vect{v}\rVert^2$, $M=\lVert\vect{v}\rVert/a$, and the
pressure-corrected engine thrust taken in the donor's own form
$T_e(p)=T_{\mathrm{vac}}-A_e p$. Integration is fixed-step RK4; a stop condition
(cutoff speed, altitude crossing, burn duration) is refined onto the condition by
bisection of the final step, so a cutoff lands on the condition rather than on
the step after it.

\subsection{Atmosphere}

US76 in two regimes. Below $86$~km geometric the seven-layer lapse-rate stack,
evaluated at geopotential altitude $H = R' h/(R'+h)$ with $R'=6356.766$~km:
\begin{equation}
T = T_b + L_b (H - H_b),
\qquad
p = \begin{cases}
p_b\left(\dfrac{T_b}{T}\right)^{g_0/(R_{\mathrm{air}}L_b)}, & L_b \neq 0,\\[2ex]
p_b\exp\!\left(-\dfrac{g_0 (H-H_b)}{R_{\mathrm{air}}T_b}\right), & L_b = 0,
\end{cases}
\qquad \rho = \frac{p}{R_{\mathrm{air}}T}.
\end{equation}
Above $86$~km the model stops being a lapse-rate one --- molecular weight varies
and the temperature profile is tabulated --- so density is log-interpolated from
the published table instead of extrapolated. That distinction is not cosmetic: an
exponential tail overstates the density at $100$~km by about $20\,\%$, and it is
precisely the near-zero that the vacuum argument above rests on.

The implementation reproduces the published density to better than $0.01\,\%$
from sea level to $120$~km (\code{rtls\_trajectory.self\_test}). Note that
$84.852$~km geopotential \emph{is} $86$~km geometric, so the two regimes meet
exactly and quoting a US76 table against the wrong one of the two shows up as a
spurious $17\,\%$ model error.

\subsection{Drag coefficient, nose-first and engines-first}
\label{subsec:cds}

The reference notebook uses $C_D=0.5$ on $4\pi=12.57\ \mathrm{m^2}$; the expert
review of this model suggested $C_D\approx1.0$. The two differ by $1.7\times$ in
$C_D S$ and the discrepancy is not a disagreement about the same number: they
describe \textbf{different attitudes}.

\begin{description}
\item[Ascent, nose-first.] Flying solo the stage carries a nose fairing over the
  open interstage, a $6.0$~m cone of $17^\circ$ half-angle on the $3.66$~m body.
  A pointed cone-cylinder at zero incidence has $C_D\approx0.30$ subsonic, a
  transonic peak of $0.55$ just past $M=1$, decaying to $0.24$ by $M=8$.
  Referenced to $S=\pi(D/2)^2=10.52\ \mathrm{m^2}$, that is a peak
  $C_D S = 5.79\ \mathrm{m^2}$.
\item[Return, engines-first.] A returning booster flies engines-first, so the
  leading end is a flat, separated base --- a bluff body, not a cone. The donor
  already carries this as a Mach table, $C_A = 0.85\ldots1.40$ on the same
  $10.52\ \mathrm{m^2}$, i.e.\ $C_D S = 8.94\ldots14.73\ \mathrm{m^2}$.
\end{description}

So the review is right about the return and the notebook's figure is closer to
the ascent value, and since the donor owns every metre below the entry interface,
the coefficient that matters for the descent is settled by the code that flies
it. This module uses the cone table for the ascent and the donor's own table for
the short stretch of its coast that is still in measurable air, so the two halves
of the flight do not disagree across the handover.

One consequence worth recording: the task's own \emph{Architecture} table quotes
the drag on an $80$~t stage at $60$~km as $0.025\ \mathrm{m/s^2}$, which is
$C_D S \approx 5.2\ \mathrm{m^2}$ --- the notebook's ascent-like value applied to
a vehicle flying backwards. With the donor's coefficient it is
$0.047\ \mathrm{m/s^2}$. Both are negligible against thrust, so the conclusion
stands, but the figure is low by $1.9\times$.

\subsection{Ascent}

The nose fairing changes the ascent drag and nothing else. The pitch programme is
vertical rise, a smoothstepped kick, then a zero-angle-of-attack gravity turn,
with a constant steering bias $\beta$ permitted only above $h_{\mathrm{free}}=40$~km
where $q<300$~Pa and steering is aerodynamically free:
\begin{equation}
\hat{\vect{u}}_T(t) =
\begin{cases}
\hat{\vect{u}}, & t<t_p,\\[1ex]
\cos(\psi\tilde s)\,\hat{\vect{u}} + \sin(\psi\tilde s)\,\hat{\vect{e}},
  \quad \tilde s = s^2(3-2s),\ s=\dfrac{t-t_p}{t_r}, & t_p \le t < t_p+t_r,\\[2ex]
\hat{\vect{v}}, & h<h_{\mathrm{free}},\\[1ex]
\cos\beta\,\hat{\vect{v}} + \sin\beta\,\hat{\vect{n}}, & h\ge h_{\mathrm{free}},
\end{cases}
\end{equation}
with $\hat{\vect{n}}\perp\hat{\vect{v}}$ on the downrange side. Throttle holds an
axial acceleration cap:
\begin{equation}
\tau = \min\!\left(1,\ \max\!\left(\tau_{\min},\ \frac{a_{\mathrm{cap}} g_0 m}{n_e T_e(p)}\right)\right).
\end{equation}

\paragraph{Why the cap is a free variable.} This is the whole difficulty of a
solo ascent. Without a $107.5$~t upper stack the stage is light, and flown at
full thrust it reaches $1579\ \mathrm{m/s}$ while still only $54$~km up ---
$6.5$~km short of the cutoff condition, with $\beta$ already saturated against
its bound and nowhere left to go. CRS-9 reached the same speed at $60.4$~km
because it was dragging the upper stack; a solo stage accelerates too well to
arrive that high that slowly. The stage therefore has to be held back to fly a
trajectory it could not otherwise reach, so $a_{\mathrm{cap}}$ is solved for
rather than fixed, and the value it settles on --- $2.79$~g --- is a result. The
burn is not much longer for it ($109.9$~s against $110.6$); it is a different
shape, climbing more steeply for the same speed.

\paragraph{The solve.} Four unknowns
$(m_0,\ t_p,\ \psi,\ \beta,\ a_{\mathrm{cap}})$ --- five, counting the cap ---
against three conditions evaluated at the instant $\lVert\vect{v}\rVert$ reaches
the cutoff speed: altitude, ground-track range, and the mass remaining. The mass
condition is supplied by the caller and is what makes the return budget close.
A fourth residual takes up the remaining freedom: either a small penalty on
$\beta$, so the solve prefers the least-steering trajectory, or a target flight
path angle when one is asked for. Solved by \code{scipy} trust-region least
squares on a coarse step, then re-flown at the fine step.

\paragraph{The flight path angle is not pinned by the cutoff spec.} Altitude,
speed and downrange are three numbers; the direction of the velocity is a fourth
and they do not determine it. It matters, because apogee follows from the
vertical component almost alone. Sweeping it shows the trade: the least-steering
solution cuts off at $\gamma\approx50^\circ$, costs $151$~t of propellant in
total and reaches $168$~km; forcing $\gamma$ down towards the $30$--$40^\circ$
the task assumes costs $171$~t and reaches $153$~km, and below about $42^\circ$
the steering bias saturates and the trajectory is not reachable at all under a
zero-incidence gravity turn through the atmosphere. CRS-9's measured $128.7$~km
apogee would need $\gamma\approx32^\circ$, which this ascent cannot produce while
also arriving at $60.4$~km only $35$~km downrange.

\subsection{Turnaround and boostback}

\paragraph{The turnaround is flown by the engines, not by the cold gas.} Out of
the atmosphere the grid fins produce no torque and, with the engines off, neither
does the gimbal, so the obvious model gives the whole flip to the RCS. Two
$500$~N nozzles at the $40.5$~m station against $J\approx6.2\times10^6\
\mathrm{kg\,m^2}$ give $\alpha_{\mathrm{RCS}}=0.31\ \mathrm{deg/s^2}$, and a
bang-bang slew of $150^\circ$ then takes $44$~s --- during which the stage
travels a further $49$~km downrange, putting maximum downrange at $97$~km. It is
also slower than the real vehicle, whose MECO-to-boostback gap on CRS-9 was
$37$~s \emph{including} stage separation.

What actually happens is that the stage settles its propellant for
$t_s=8$~s, relights, and lets the gimbal finish the turn:
\begin{equation}
\alpha_{\mathrm{lit}}
 = \frac{n_r F_r\lvert z_r - z_{\mathrm{com}}\rvert + T\sin\delta_{\max}\,z_{\mathrm{com}}}{J}
 \approx 15.7\ \mathrm{deg/s^2},
\end{equation}
fifty times the cold gas, so the same $150^\circ$ takes $5.7$~s. The burn and the
turn are one manoeuvre, so the thrust direction slews \emph{during} the burn,
bang-bang within $\alpha_{\mathrm{lit}}$, and the solver integrates the fact that
the thrust genuinely points the wrong way for the first few seconds rather than
assuming it away. This one modelling choice moves maximum downrange from $97$~km
to $61$~km, i.e.\ from outside the $59$--$70$~km band to inside it.

\paragraph{The boostback as a boundary-value problem.} This is a targeting
problem, not a landing problem, which is why the SCvx formulation of
Section~\ref{sec:guidance} does not apply: its terminal condition is ``at the
origin with near-zero velocity'' and the stage at boostback cutoff is $100$~km up
and moving at a kilometre a second. What is solved is the return leg of the
reference notebook's phase~2, posed on the cutoff state:

\begin{quote}
given the state at engine cutoff, find a burn --- attitude $\chi$ held fixed in
the inertial frame, duration $t_b$ --- whose ballistic continuation reaches the
launch site.
\end{quote}

That is one equality constraint against two unknowns, so there is a
one-parameter family of solutions and the spare freedom buys the cheapest of
them:
\begin{equation}
\min_{\chi}\ t_b(\chi)
\quad\text{subject to}\quad
d_{\mathrm{impact}}\big(\chi, t_b\big) = d_{\mathrm{aim}} ,
\end{equation}
which at fixed engine count and throttle is minimum propellant. For each $\chi$
the duration is found by bisection --- the miss is monotone in $t_b$, since a
longer burn always leaves the stage further up-range --- and $\chi$ is then
optimised by bounded scalar minimisation. Apogee is left alone by this, so it
remains an independent check rather than something the solve was told to hit.

\paragraph{The objective is degenerate, and needs one constraint.} Scanning
$\chi$ at a fixed cutoff state, the propellant varies by $2\,\%$ ---
$32.64$~t at $172.5^\circ$ against $33.29$~t at $185^\circ$ --- while apogee
moves from $163$~km to $132$~km and the vertical velocity at the entry interface
from $-1329$ to $-1090\ \mathrm{m/s}$. A minimum-fuel objective cannot
discriminate between those, and left unconstrained it drifts below $180^\circ$
into thrusting \emph{upward}, lofting the stage past $200$~km to save thirty
kilograms and handing the donor a stage arriving $500\ \mathrm{m/s}$ faster. The
objective cannot see that cost because it stops at boostback cutoff, while the
bill arrives later in the entry burn.

The attitude is therefore constrained to $\chi\ge180^\circ$: the boostback
thrust may not have an upward component. This is a statement about what a
boostback \emph{is} --- a retrograde manoeuvre --- rather than a tuning choice,
it costs $0.1\,\%$ of the propellant, and apogee still falls out of the solve
rather than being targeted.

\paragraph{The bracket is the propellant, not a round number.} The burn-time
bracket must end just short of a dry tank. A fixed upper bound of $90$~s asks
three Merlins for $82$~t from a stage carrying $61$, so the search integrates
negative propellant mass, where $T/m$ diverges through zero. The resulting
trajectories are meaningless rather than merely inaccurate, and a meaningless
trajectory can still bracket a root: with the round bound, four attitudes in the
middle of the useful range reported themselves unreachable. The equations of
motion additionally refuse to thrust a stage at or below $m_{\mathrm{dry}}$.

\subsection{Handover and the aim-point loop}
\label{subsec:handover}

The arc is targeted on its \emph{vacuum} impact point, which is the honest
calculation above the interface but is not where the stage lands: the real descent has
an atmosphere and the arc falls short. Rather than model the donor's drag a
second time, the aim point is corrected against the donor itself --- fly it,
measure the touchdown miss $\epsilon$, set $d_{\mathrm{aim}} \leftarrow
d_{\mathrm{aim}} - 0.9\,\epsilon$, fly it again. The relationship is very nearly
one-to-one, so two or three passes suffice. This is a shooting method with the
donor as the plant.

The handover itself is a state vector through the \code{custom} scenario added to
the donor (Section~\ref{sec:donorpatch}): altitude, ground-track range, both
velocity components, propellant remaining, and an attitude built along
$-\hat{\vect{v}}$ rather than near-vertical, because a stage arriving from a
boostback is flying engines-first at some $9^\circ$ off vertical, well outside
the $\pm2^\circ$ every shipped scenario assumes.

\paragraph{Why the interface sits at 65\,km.} The donor's entry burn is armed
and cut on a \emph{predicted peak dynamic pressure}, and at $70$~km, where
$q=62$~Pa, that prediction falls below the cut threshold while about $20$~t is
still aboard. The burn stops early, which leaves a $55$~t stage for a landing
burn flown on one engine. That nets
\begin{equation}
a_{\mathrm{net}} = \frac{T_{\mathrm{SL}}}{m} - g_0
 = \frac{845}{55.6} - 9.81 = 5.4\ \mathrm{m/s^2},
\qquad
\frac{v^2}{2a_{\mathrm{net}}} = \frac{380^2}{10.8} = 13.4\ \mathrm{km},
\end{equation}
and $13$~km is exactly where the landing burn then ignites; it runs the tank dry
on the way down. Sweeping the handover propellant at that interface, every load
from $26$ to $36$~t crashes, and more propellant is monotonically worse.

Below $66$~km the entry burn runs instead to its $7$~t fuel floor, leaving a
$33$~t stage whose one-engine burn needs only $2.7$~km --- the regime the donor
was built and tested in --- and every interface altitude from $66$~km down lands
at $2.6$ to $4.3\ \mathrm{m/s}$. The interface is therefore placed at $65$~km,
clear of the cliff between $68$ and $66$.

This costs nothing in validity. The vacuum argument of the specification is made
at $70$~km, but it is an argument about dynamic pressure rather than about a
particular altitude: at $65$~km, $q\approx130$~Pa and the drag deceleration on
the returning stage is $0.03\ \mathrm{m/s^2}$, a few hundredths of a percent of
the entry burn's thrust acceleration and two orders of magnitude below where the
approximation begins to matter.

\subsection{The one donor patch}
\label{sec:donorpatch}

Every shipped scenario starts with essentially zero horizontal velocity:
\code{scenario\_init} gives $V_X$ and $V_Y$ a $\pm5\ \mathrm{m/s}$ dispersion and
nothing else, so a stage still flying back towards the pad at a couple of hundred
metres a second could not be expressed at all. The patch adds a
\code{vx0} field to \code{ScenDef}, a \code{custom} scenario whose initial
conditions come from the command line, and a retrograde initial attitude. Every
shipped scenario leaves \code{vx0} at zero and is bit-identical as a result;
\code{--selftest} passes and the \code{entry} scenario reproduces its documented
touchdown exactly ($3.19\ \mathrm{m/s}$, $24.07$~m).

\section{Projection for display}
\label{sec:projection}

\subsection{Vertical plane}

The downrange axis is the horizontal direction of the initial position, and the
cross-range axis completes the frame:
\begin{equation}
\eh = \frac{\vect{r}_{0,xy}}{\lVert\vect{r}_{0,xy}\rVert},
\qquad
\ec = \hat{\vect{z}}\times\eh .
\end{equation}
Positions and velocities project as $d=\vect{r}\cdot\eh$,
$c=\vect{r}\cdot\ec$, with altitude $r_z$ taken directly. The body and thrust
directions project the same way after rotation into the inertial frame:
\begin{equation}
\hat{\vect{b}}_I = \Rot(\quat{q})\vect{e}_3,
\qquad
\hat{\vect{b}}_{\text{2D}} = \left(\hat{\vect{b}}_I\cdot\eh,\ \hat b_{I,z}\right).
\end{equation}

\subsection{Derived angles}

\begin{align}
\theta_{\mathrm{tilt}} &= \arccos\bigl(\mathrm{clip}(1-2(q_1^2+q_2^2),-1,1)\bigr), \\
\delta_{\mathrm{gimbal}} &= \operatorname{atan2}\!\left(\lVert(u_x,u_y)\rVert,\ \max(u_z,\epsilon)\right), \\
\alpha_{\mathrm{AoA}} &= \arccos\!\left(\mathrm{clip}\!\left(
  \frac{-\vect{v}}{\lVert\vect{v}\rVert}\cdot\hat{\vect{b}}_I,-1,1\right)\right),
  \quad \text{undefined for } \lVert\vect{v}\rVert<1\ \mathrm{m/s}.
\end{align}

\subsection{Drawing convention}

The hull is drawn with its nozzle on the logged point rather than its centre of
mass on it. For \code{pdg} data this is a deliberate offset of the engine arm
$\ell$, because that model's state is the centre of mass and its touchdown
condition is the centre of mass reaching $z=0$; drawing the hull around it would
bury the vehicle in the deck. For donor data the logged height is already the
vehicle base, so the same rule is exact. Every plotted number remains the raw
logged value.

\subsection{Camera}

Included for completeness; it affects presentation only. The view is a square
window at unit aspect throughout. On a flight that \emph{reaches} above
$h_{\mathrm{zoom}}$ the camera holds the whole trajectory, then transitions to a
chase window of width $w_c$ over a fixed wall-clock interval during which the
replay is held. The trigger runs on a running minimum of altitude taken from the
last sample above $h_{\mathrm{zoom}}$, not from the first sample of the flight:
every profile before the RTLS one began at its highest point and came down, so
testing the first sample was the same thing as testing the highest, and a
running minimum from $t=0$ was the same thing as one from apogee. A flight that
lifts off from the pad breaks both --- it starts at zero altitude, which switched
the camera off, and latched a zero into the running minimum before the stage had
left the ground.
The transition interpolates geometrically in scale and smoothly in time:
\begin{equation}
\tilde s = s^2(3-2s),
\qquad
w(s) = W\left(\frac{w_c}{W}\right)^{\tilde s},
\qquad
\vect{p}(s) = (1-\tilde s)\,\vect{p}_{\mathrm{scene}} + \tilde s\,\vect{p}_{\mathrm{booster}} .
\end{equation}
The scale bar takes the largest value of the form $\{1,2,5\}\times10^n$ not
exceeding $w/4$.

\section{Assumptions and open questions}
\label{sec:limits}

Ordered roughly by how much they could matter.

\begin{enumerate}
\item \textbf{Constant inertia in the pdg plant.} $\vect{J}$ in
  \eqref{eq:wdot} is fixed while the vehicle burns a third of its mass. The
  donor reconstruction computes a mass-varying inertia; the C\texttt{++} plant
  does not. The attitude response therefore stiffens less than it should late in
  the burn.

\item \textbf{In-plane assumption for donor data.} Both the range unfolding and
  the lean reconstruction assume all horizontal motion lies in one vertical
  plane. The log carries no azimuth, so this cannot be verified. Any real
  azimuth change appears as an in-plane reversal, which the slew limiter then
  renders as a physically-paced pitch-over that did not necessarily occur.

\item \textbf{Lean sign from lateral acceleration.} The inference ignores
  aerodynamic side force, which during the aero descent is not small. The sign is
  held through unpowered stretches, so a genuine reversal there is missed until
  the next burn.

\item \textbf{Autopilot gimbal saturation exceeds the guidance envelope.} The
  lateral clamp at $\tfrac{1}{2}\lVert\vect{f}\rVert$ permits $30^\circ$ against
  a guidance constraint of $12^\circ$.

\item \textbf{Linearised lower thrust bound.} Unlike the 3-DoF lossless
  relaxation, the 6-DoF minimum-thrust constraint is linearised about the
  reference, so the feasible set depends on the iterate.

\item \textbf{Grid-fin authority is an upper bound.} Equation
  \eqref{eq:alphamax} uses maximum deflection at the instantaneous dynamic
  pressure with a constant $C_{N\alpha}$, no stall, no Mach dependence and no
  blanketing. It is generous; the $\omega_\perp$ cap is what actually binds.

\item \textbf{Isp pressure blend.} An isothermal exponential rather than the
  US76 pressure profile the donor itself integrates.

\item \textbf{No contact model in pdg.} Touchdown is a terminal condition, not a
  landing: legs, crush and tip-over exist only in the donor.

\item \textbf{Roll is uncontrolled and unmodelled in the projection.} A single
  gimballed engine has no roll authority in the pdg model, and the 2-D
  projection discards roll entirely.

\item \textbf{No Earth rotation and no cross-range in the RTLS model.} A due-east
  launch from $28.5^\circ$ starts with $408\ \mathrm{m/s}$ of eastward ground
  speed and picks up a Coriolis cross-range. Section~\ref{sec:rtls} is planar and
  non-rotating; so is the donor's frame, so adding it would have to stop at the
  handover anyway. It biases the boostback $\Delta v$ and the ground track, and
  the aim-point loop absorbs the residual rather than accounting for it.

\item \textbf{The RTLS ascent has no attitude dynamics.} Section~\ref{sec:rtls}
  commands a thrust direction and the vehicle takes it. The slew during the
  boostback is rate-limited by the real gimbal authority, and the coast turn by
  the real RCS authority, but neither is a closed attitude loop with an actuator
  model, and there is no check that the commanded pitch rate is trackable
  through the ascent's own dynamic pressure.

\item \textbf{The return is wind-limited and the nominal flight is flown in
  light wind.} Under the donor \code{entry} scenario's $8\ \mathrm{m/s}$ surface
  wind and severe Dryden turbulence, the same trajectory lands at
  $2.8\ \mathrm{m/s}$ or crashes at $27.6\ \mathrm{m/s}$ depending only on the
  handover propellant load. This is the architectural fragility of a stage that
  cannot hover --- one chance at the ignition point --- and it is not fixed by
  anything in this model. The nominal RTLS flight is therefore flown at
  $3\ \mathrm{m/s}$ and moderate turbulence, and the metadata records it.
\end{enumerate}

\section{Parameters}

\subsection{pdg vehicle, default case}

\begin{center}
\begin{tabular}{llr}
\toprule
symbol & meaning & value \\
\midrule
$m_{\mathrm{wet}}$ & wet mass & $15\,000$ kg \\
$m_{\mathrm{dry}}$ & dry mass & $10\,000$ kg \\
$I_{sp}$ & specific impulse & $300$ s \\
$T_{\min}$, $T_{\max}$ & thrust bounds & $80$, $250$ kN \\
$\vect{J}$ & inertia (constant) & $(8, 8, 1)\times10^5\ \mathrm{kg\,m^2}$ \\
$\vect{r}_T$ & thrust point & $(0,0,-12)$ m \\
$\theta_{\max}$ & tilt limit & $70^\circ$ \\
$\delta_{\max}$ & gimbal limit & $12^\circ$ \\
$\gamma_{gs}$ & glide slope & $20^\circ$ \\
$\omega_{\max}$ & body rate limit & $60^\circ$/s \\
$K$ & SCvx nodes & $21$ (default) \\
$w_\nu$ & virtual control weight & $10^3$ \\
\bottomrule
\end{tabular}
\end{center}

\subsection{Donor vehicle constants used by the reconstruction}

Taken from \code{external/booster-lander-simulator/core/constants.h}.

\begin{center}
\begin{tabular}{llr}
\toprule
symbol & meaning & value \\
\midrule
$L$ & vehicle length & $47.7$ m \\
$R$ & radius & $1.83$ m \\
$m_{\mathrm{dry}}$ & dry mass incl.\ legs and fins & $25\,600$ kg \\
$z_{\mathrm{dry}}$ & dry centre of mass above base & $12.4$ m \\
$A$ & tank column cross-section & $9.9\ \mathrm{m^2}$ \\
$z_{\mathrm{base}}$ & column base & $1.6$ m \\
$\rho$ & mixture density (assumed, $O/F=2.56$) & $1048\ \mathrm{kg/m^3}$ \\
$T_{\mathrm{sl}}$ & engine thrust, sea level & $845$ kN \\
$I_{sp}^{\mathrm{sl}}$, $I_{sp}^{\mathrm{vac}}$ & specific impulse & $282$, $311$ s \\
$\delta_{\max}$ & gimbal limit & $5^\circ$ \\
$n_f$, $S_f$, $C_{N\alpha}$, $\delta_f$ & fin pair, area, slope, deflection &
  $2$, $2.4\ \mathrm{m^2}$, $3.0$/rad, $20^\circ$ \\
$z_f$ & fin station & $45.0$ m \\
$n_r$, $F_r$, $z_r$ & RCS couple, thrust, station & $2$, $500$ N, $40.5$ m \\
\bottomrule
\end{tabular}
\end{center}

\section{Source map}

\begin{center}
\begin{tabular}{ll}
\toprule
section & implementation \\
\midrule
\ref{sec:vehicle} Vehicle model & \code{external/pdg/src/rocket.cpp} \\
\ref{sec:guidance} Guidance & \code{external/pdg/src/scvx.cpp}, \code{lcvx.cpp}, \code{socp.cpp} \\
\ref{sec:closedloop} Plant, autopilot & \code{external/pdg/src/sim.cpp} \\
\ref{sec:closedloop} Feasibility screen & \code{bridge/pdg\_flight.cpp} \\
\ref{sec:donor} Donor reconstruction & \code{bridge/bls\_flight.py} \\
\ref{sec:projection} Projection, geometry & \code{viz/flight\_data.py} \\
\ref{sec:projection} Camera & \code{viz/land2d\_mpl.py} \\
\bottomrule
\end{tabular}
\end{center}

\end{document}
